\section{Conclusion}

We present a framework that connects target-space definition, geometric quality measurement and paired-corner structural consensus to the study of human layout annotation. $Q$ captures discrepancies in coverage, boundary position, direction and height, while consensus combines distinct-person support and upper--lower observations into a connected layout. Doorway-boundary viewpoints extend the annotation problem to camera positions where adjoining enclosures require an explicit target and closure boundary. Together, these contributions make differences in both the intended space and its geometric expression available for analysis and review.

The empirical findings show how people shape collective geometry. Early team expansion improves mean quality, while later additions chiefly improve the lower tail and consistency within the observed pools. Class composition supplies an average trend; specific incoming and retained members determine which turns survive and where boundaries fall. Across room views, structural habits recur alongside changes in local corner choices. These findings guide team organization and targeted re-annotation: repeatedly supported structures can be checked against the reference, and disputed enclosures can be resolved through explicit target rules and expert review. The framework provides a basis for producing higher-quality layout datasets for model training and evaluation.
